% Options for packages loaded elsewhere
\PassOptionsToPackage{unicode}{hyperref}
\PassOptionsToPackage{hyphens}{url}
\PassOptionsToPackage{dvipsnames,svgnames,x11names}{xcolor}
\documentclass[
  11pt,
]{article}
\usepackage{xcolor}
\usepackage[margin=1in]{geometry}
\usepackage{amsmath,amssymb}
\setcounter{secnumdepth}{-\maxdimen} % remove section numbering
\usepackage[T1]{fontenc}
\usepackage[utf8]{inputenc}
\usepackage{textcomp} % provide euro and other symbols
\usepackage{lmodern}
\usepackage{helvet} % built-in Helvetica clone (no xelatex needed)
\renewcommand{\familydefault}{\sfdefault} % use sans-serif as the body font
% Use upquote if available, for straight quotes in verbatim environments
\IfFileExists{upquote.sty}{\usepackage{upquote}}{}
\IfFileExists{microtype.sty}{% use microtype if available
  \usepackage[]{microtype}
  \UseMicrotypeSet[protrusion]{basicmath} % disable protrusion for tt fonts
}{}
\makeatletter
\@ifundefined{KOMAClassName}{% if non-KOMA class
  \IfFileExists{parskip.sty}{%
    \usepackage{parskip}
  }{% else
    \setlength{\parindent}{0pt}
    \setlength{\parskip}{6pt plus 2pt minus 1pt}}
}{% if KOMA class
  \KOMAoptions{parskip=half}}
\makeatother
\usepackage{color}
\usepackage{fancyvrb}
\newcommand{\VerbBar}{|}
\newcommand{\VERB}{\Verb[commandchars=\\\{\}]}
\DefineVerbatimEnvironment{Highlighting}{Verbatim}{commandchars=\\\{\}}
% Add ',fontsize=\small' for more characters per line
\newenvironment{Shaded}{}{}
\newcommand{\AlertTok}[1]{\textcolor[rgb]{1.00,0.00,0.00}{\textbf{#1}}}
\newcommand{\AnnotationTok}[1]{\textcolor[rgb]{0.38,0.63,0.69}{\textbf{\textit{#1}}}}
\newcommand{\AttributeTok}[1]{\textcolor[rgb]{0.49,0.56,0.16}{#1}}
\newcommand{\BaseNTok}[1]{\textcolor[rgb]{0.25,0.63,0.44}{#1}}
\newcommand{\BuiltInTok}[1]{\textcolor[rgb]{0.00,0.50,0.00}{#1}}
\newcommand{\CharTok}[1]{\textcolor[rgb]{0.25,0.44,0.63}{#1}}
\newcommand{\CommentTok}[1]{\textcolor[rgb]{0.38,0.63,0.69}{\textit{#1}}}
\newcommand{\CommentVarTok}[1]{\textcolor[rgb]{0.38,0.63,0.69}{\textbf{\textit{#1}}}}
\newcommand{\ConstantTok}[1]{\textcolor[rgb]{0.53,0.00,0.00}{#1}}
\newcommand{\ControlFlowTok}[1]{\textcolor[rgb]{0.00,0.44,0.13}{\textbf{#1}}}
\newcommand{\DataTypeTok}[1]{\textcolor[rgb]{0.56,0.13,0.00}{#1}}
\newcommand{\DecValTok}[1]{\textcolor[rgb]{0.25,0.63,0.44}{#1}}
\newcommand{\DocumentationTok}[1]{\textcolor[rgb]{0.73,0.13,0.13}{\textit{#1}}}
\newcommand{\ErrorTok}[1]{\textcolor[rgb]{1.00,0.00,0.00}{\textbf{#1}}}
\newcommand{\ExtensionTok}[1]{#1}
\newcommand{\FloatTok}[1]{\textcolor[rgb]{0.25,0.63,0.44}{#1}}
\newcommand{\FunctionTok}[1]{\textcolor[rgb]{0.02,0.16,0.49}{#1}}
\newcommand{\ImportTok}[1]{\textcolor[rgb]{0.00,0.50,0.00}{\textbf{#1}}}
\newcommand{\InformationTok}[1]{\textcolor[rgb]{0.38,0.63,0.69}{\textbf{\textit{#1}}}}
\newcommand{\KeywordTok}[1]{\textcolor[rgb]{0.00,0.44,0.13}{\textbf{#1}}}
\newcommand{\NormalTok}[1]{#1}
\newcommand{\OperatorTok}[1]{\textcolor[rgb]{0.40,0.40,0.40}{#1}}
\newcommand{\OtherTok}[1]{\textcolor[rgb]{0.00,0.44,0.13}{#1}}
\newcommand{\PreprocessorTok}[1]{\textcolor[rgb]{0.74,0.48,0.00}{#1}}
\newcommand{\RegionMarkerTok}[1]{#1}
\newcommand{\SpecialCharTok}[1]{\textcolor[rgb]{0.25,0.44,0.63}{#1}}
\newcommand{\SpecialStringTok}[1]{\textcolor[rgb]{0.73,0.40,0.53}{#1}}
\newcommand{\StringTok}[1]{\textcolor[rgb]{0.25,0.44,0.63}{#1}}
\newcommand{\VariableTok}[1]{\textcolor[rgb]{0.10,0.09,0.49}{#1}}
\newcommand{\VerbatimStringTok}[1]{\textcolor[rgb]{0.25,0.44,0.63}{#1}}
\newcommand{\WarningTok}[1]{\textcolor[rgb]{0.38,0.63,0.69}{\textbf{\textit{#1}}}}
\usepackage{longtable,booktabs,array}
\usepackage{calc} % for calculating minipage widths
% Correct order of tables after \paragraph or \subparagraph
\usepackage{etoolbox}
\makeatletter
\patchcmd\longtable{\par}{\if@noskipsec\mbox{}\fi\par}{}{}
\makeatother
% Allow footnotes in longtable head/foot
\IfFileExists{footnotehyper.sty}{\usepackage{footnotehyper}}{\usepackage{footnote}}
\makesavenoteenv{longtable}
\setlength{\emergencystretch}{3em} % prevent overfull lines
\providecommand{\tightlist}{%
  \setlength{\itemsep}{0pt}\setlength{\parskip}{0pt}}
\usepackage{newunicodechar}
\newunicodechar{→}{\textrightarrow}
\usepackage{fvextra}
\DefineVerbatimEnvironment{Highlighting}{Verbatim}{breaklines,breaknonspaceingroup,commandchars=\\\{\}}
\setlength{\emergencystretch}{3em}
\usepackage{etoolbox}
\AtBeginEnvironment{longtable}{\footnotesize}
\usepackage[most]{tcolorbox}
\renewenvironment{Shaded}{\begin{tcolorbox}[breakable,boxrule=0.4pt,colback=gray!8,colframe=gray!50,arc=2pt,left=5pt,right=5pt,top=4pt,bottom=4pt]}{\end{tcolorbox}}
\usepackage{bookmark}
\IfFileExists{xurl.sty}{\usepackage{xurl}}{} % add URL line breaks if available
\urlstyle{same}
\hypersetup{
  colorlinks=true,
  linkcolor={blue},
  filecolor={Maroon},
  citecolor={Blue},
  urlcolor={blue},
  pdfcreator={LaTeX via pandoc}}

\author{}
\date{}

\begin{document}

\section{Pacific Maritime Emissions ---
Documentation}\label{pacific-maritime-emissions-documentation}

\subsection{1. Overview}\label{overview}

This project estimates greenhouse-gas and air-pollutant emissions from
ships operating in the Exclusive Economic Zones (EEZs) of the Pacific
Island Countries (PICs), following the bottom-up activity-based method of
the \href{https://www.imo.org/en/ourwork/Environment/Pages/Fourth-IMO-Greenhouse-Gas-Study-2020.aspx}{Fourth
IMO GHG Study 2020}. Starting from Automatic Identification System (AIS)
vessel positions on the UN Global Platform, it builds an hourly
vessel-activity panel, attaches each vessel's technical specifications
from the IHS ship register, and reconstructs the engine load and
operating state of every vessel-hour. Emissions are then computed two
ways --- a full activity-based calculation for vessels matched to the
register, and a median-per-type model for the rest --- flagged for
fishing activity using Global Fishing Watch, and rolled up into monthly
country and vessel-type statistics that feed the dashboard. The result is
a per-vessel-hour emissions record covering CO\textsubscript{2}, SOx,
NOx, CH\textsubscript{4}, CO, N\textsubscript{2}O, NMVOC, PM10, PM2.5,
and black carbon.

\subsection{2. Data Sources}\label{data-sources}

\begin{itemize}
\tightlist
\item
  \textbf{AIS (UNGP):} raw satellite/terrestrial vessel positions, read
  on the UN Global Platform through the \texttt{ais} library
  (\texttt{af.get\_ais}). Provides position, speed, draught, and
  navigational status (notebook~A).
\item
  \textbf{IHS ship register (UNGP):} vessel identity and technical
  specifications, read via \texttt{af.read\_ihs\_table} ---
  \texttt{ShipData.CSV}, \texttt{tblCapacities.CSV},
  \texttt{tblBoilers.CSV}, and \texttt{tblNameHistory.CSV}. Used both to
  identify vessels (notebook~A) and to derive engine, fuel, and emission
  parameters (notebook~B). The IHS version is pinned per run
  (\texttt{ihs\_version}, e.g.\ \texttt{20260202}).
\item
  \textbf{IMF port boundaries:} port-area polygons
  (\texttt{imf\_port\_boundary.pkl}), used to flag positions inside a
  port (notebook~A).
\item
  \textbf{EEZ and coastal boundaries:} each country's EEZ polygon and a
  10\,km coastal buffer (\texttt{geoboundary2.pkl}), used to assign
  positions to a country and to the coast/at-sea zones (notebook~A).
\item
  \textbf{Global Fishing Watch (effort):} daily fishing effort by MMSI
  (\texttt{GFW\_effort\_event/\{year\}.csv}), used to flag fishing
  vessels and, with gear type, to break out fishing emissions
  (notebooks~C, D, E).
\item
  \textbf{Auxiliary engine and boiler power factors:} tabulated factors
  from the Third/Fourth IMO studies, pulled as JSON from
  \href{https://www.gabrielfuentes.org/tables_emissions/ae_ab_power.json}{gabrielfuentes.org}
  (notebook~B).
\item
  \textbf{Method and emission factors:} the Fourth IMO GHG Study 2020,
  which defines the vessel classes, correction factors, specific fuel
  consumption baselines, and per-pollutant emission factors.
\end{itemize}

\subsection{3. Process Flow}\label{process-flow}

The pipeline is a series of notebooks (A--E), each consuming the output
of the previous stage. All stages run inside the UN Global Platform.

\begin{enumerate}
\item
  \textbf{Build the hourly AIS panel (A).} Raw AIS is restricted to the
  Pacific bounding box and the PIC EEZs, matched to the IHS register to
  recover a stable \texttt{vessel\_id}, resampled to one record per
  vessel per hour, and tagged with EEZ, coast, and port flags via Sedona
  spatial joins. Output: \texttt{hourly\_ais\_v2/}.
\item
  \textbf{Prepare vessel specifications (B).} IHS \texttt{ShipData} (with
  capacities and boilers) is cleaned and gap-filled --- beam, draught,
  length, capacity/size bin, speed, power, engine RPM, engine type, and
  fuel type --- then turned into the correction factors, specific
  fuel-consumption baselines, and pollutant emission factors needed for
  the calculation. Output: \texttt{ihs\_emissions\_specs\ \{version\}}.
\item
  \textbf{Compute emissions, with IHS (C).} The hourly panel is joined to
  the vessel specs; engine load, power demand, operational phase, and
  low-load factors are derived; and per-hour emissions of every pollutant
  (main engine, pilot fuel, auxiliary engine, boiler) are produced.
  Output: \texttt{emissions\_data/with\_ihs/}.
\item
  \textbf{Compute emissions, without IHS (D).} For vessels not matched to
  the register, a median per-hour emission by AIS vessel group and
  operational phase is learned from the with-IHS output, then applied to
  the unmatched vessel-hours. Output: \texttt{emissions\_data/no\_ihs/}.
\item
  \textbf{Aggregate statistics (E).} The two emission sets are summed by
  month, country, vessel type, operational phase, fishing flag, and gear;
  CO\textsubscript{2}-equivalent totals are computed; and the dashboard
  tables are written. Output: \texttt{emissions\_data/stats/}.
\end{enumerate}

\subsection{4. Assumptions}\label{assumptions}

\subsubsection{Vessel identity (AIS to
IHS)}\label{vessel-identity}

Each AIS vessel is matched to the IHS register on exact IMO/MMSI/name
criteria, with a fuzzy vessel-name fallback and a name-history pass for
the remainder. A matched vessel takes the identity
\texttt{imo-\textless{}LRIMOShipNo\textgreater{}}; an unmatched vessel
keeps \texttt{mmsi-\textless{}mmsi\textgreater{}} and is routed to the
without-IHS calculation rather than being assigned a guessed identity.
The IHS match type is retained for auditing.

\subsubsection{Operational phase}\label{operational-phase}

Each vessel-hour is assigned an operating phase from its speed over
ground (SOG) and its proximity flags, which then selects the
auxiliary/boiler power and low-load factors:

\begin{longtable}[]{@{}
  >{\raggedright\arraybackslash}p{0.45\textwidth}
  >{\raggedright\arraybackslash}p{0.25\textwidth}@{}}
\toprule
\textbf{Condition} & \textbf{Phase} \\
\midrule
\endhead
SOG $\leq$ 1 kn and inside a port & At Berth \\
SOG $\leq$ 3 kn & Anchored \\
3 \textless{} SOG $\leq$ 5 kn, near port or coast & Manoeuvring \\
otherwise & Sea \\
\bottomrule
\end{longtable}

\subsubsection{Filling missing vessel
specifications}\label{filling-missing-specs}

The vessel specifications used in the calculation (beam, draught, length,
capacity, speed, power, engine RPM) are filled in a fixed order: use the
IHS reference field where present; otherwise estimate it with a
single-factor or IMO-style regression on related fields; otherwise fall
back to the median for the vessel's class and size bin. A value of 0 in
these fields is treated as missing. Where the WB method departs from the
IMO demo, it adds further IHS reference fields to the regression ---
\texttt{BreadthMoulded} for beam, \texttt{LengthRegistered} and
\texttt{LengthBetweenPerpendicularsLBP} for length, and
\texttt{Speed}/\texttt{Speedservice} for max speed. Engine RPM is used
only as a band (slow / medium / high speed), so the band matters rather
than the exact value. The source of each filled value is recorded in a
companion \texttt{\_basis} column.

\subsubsection{Deviations from the IMO
method}\label{imo-deviations}

Two filling rules differ deliberately from the Fourth IMO GHG Study:

\begin{itemize}
\tightlist
\item
  \textbf{Missing capacity.} Vessels whose capacity is unknown are placed
  in a separate size bin rather than imputed to the median size, then
  assigned the modal size bin of their vessel class. This avoids giving
  unknown-capacity vessels a spuriously precise size.
\item
  \textbf{Engine type and fuel type.} These are filled \emph{jointly}
  (most common consistent combination per class and size bin) instead of
  independently, so that physically inconsistent pairings --- e.g.\ an
  oil engine (SSD/MSD/HSD) reported with LNG fuel --- are not produced.
  The IMO assumption that little fuel-type data is missing does not hold
  well in this dataset.
\end{itemize}

\subsubsection{Engine, fuel, and emission
factors}\label{engine-fuel-factors}

\begin{itemize}
\tightlist
\item
  \textbf{Specific fuel consumption (SFC).} Baseline SFC is assigned by
  engine type, fuel, and engine generation, where generation is a
  build-year tier split at 1984 and 2001.
\item
  \textbf{Auxiliary engine and boiler.} Both are assumed to burn the same
  fuel as the main engine; their power demand follows the Third IMO GHG
  Study factors by vessel type and size.
\item
  \textbf{Pilot fuel.} LNG engines (LNG-Otto, LNG-Diesel) are assumed to
  burn an MDO pilot fuel.
\item
  \textbf{Fuel sulphur content.} Assumed at 2.6\% for HFO and 0.07\% for
  MDO; this drives the SOx and PM factors.
\item
  \textbf{CO\textsubscript{2} emission factors} (tonnes CO\textsubscript{2}
  per tonne fuel): HFO 3.114, MDO 3.206, LNG 2.750, Methanol 1.375.
\end{itemize}

\subsubsection{NOx tier and Emission Control
Areas}\label{nox-tier}

NOx factors are assigned by engine RPM band and IMO NOx Tier, with the
tier set from build year (Tier I from 2000, Tier II from 2011, Tier III
from 2016). The tier values used assume operation \emph{outside} an
Emission Control Area, so the stricter Tier III limits are not applied
--- a conservative choice given that the Pacific region is not an ECA.

\subsubsection{Power correction
factors}\label{correction-factors}

Main-engine power demand is adjusted by a weather correction factor, a
constant hull-fouling factor (0.917), and a speed correction factor that
is applied only to large container ships and cruise ships (left at 1 for
all other vessels).

\subsubsection{Emissions and
CO\textsubscript{2}-equivalent}\label{emissions-and-co2eq}

Per-hour emissions are scaled by the time gap to the next position
(capped at 24~hours) so that gaps in coverage do not over- or
under-count. CO\textsubscript{2}-equivalent uses GWP100 factors ---
CH\textsubscript{4} $\times$ 27 and N\textsubscript{2}O $\times$ 273 ---
added to CO2. In code:

\begin{Shaded}
\begin{Highlighting}[]
df['ch4_co2eq']   = df['_ch4_e'] * 27      # GWP100 (CH4)
df['n2o_co2_eq']  = df['_n2o_e'] * 273     # GWP100 (N2O)
df['co2eq_total'] = df['_co2_f'] + df['ch4_co2eq'] + df['n2o_co2_eq']
\end{Highlighting}
\end{Shaded}

\subsubsection{Scope and exclusions}\label{scope-and-exclusions}

Coverage is the EEZs of 14 Pacific Island Countries. Longitudes are
transposed to a 0--360\textdegree{} frame so the area stays contiguous
across the antimeridian, and the Kiribati sub-regions (Phoenix, Gilbert,
and Line groups) are reported together as Kiribati. Vessels fuelled by
ethane or LPG are excluded from the CO\textsubscript{2}/SFC factors, as
their numbers are negligible.

\subsection{5. Deployment}\label{deployment}

\subsubsection{Requirements}\label{requirements}

\begin{itemize}
\tightlist
\item
  A \textbf{UN Global Platform (UNGP) account}
  (\href{https://datalab.officialstatistics.org}{datalab.officialstatistics.org}).
  Unlike a local notebook project, \emph{every} stage here runs inside
  UNGP, because all of them use the \texttt{ais} library, the Spark
  cluster, and the AIS/IHS data held on the platform.
\item
  \textbf{PySpark with Apache Sedona}, provided by the platform's
  \texttt{ais} library, for the spatial joins and the distributed
  emissions calculation.
\item
  \textbf{Python 3} with the scientific stack (pandas, numpy, geopandas,
  scikit-learn) for the driver-side cleaning and aggregation.
\item
  \textbf{Global Fishing Watch effort files}
  (\texttt{GFW\_effort\_event/}) staged in the project storage, and
  \textbf{outbound internet} for the auxiliary/boiler power-factor JSON.
\end{itemize}

\subsubsection{Platform notes}\label{platform-notes}

The notebooks run on UNGP's current environment (Onyxia / EMR on EKS with
Spark Connect, Spark~3.5). Each notebook creates its session with
\texttt{SparkSession.builder.getOrCreate()} (which reads
\texttt{SPARK\_REMOTE}), reads AIS and the IHS register through the
\texttt{ais} library, and writes all project data to the user's personal
bucket under \texttt{AWS\_WORKING\_DIRECTORY\_PATH}. The shared AIS data
bucket is read-only.

\subsubsection{Execution}\label{execution}

The notebooks run in order, A through E; each reads the previous stage's
output and writes the input for the next. Every notebook opens with a
configuration cell that holds the user inputs in one place --- storage
paths, the IHS version, and the month ranges to process.

\subsection{6. Files}\label{files}

\begin{longtable}[]{@{}
  >{\raggedright\arraybackslash}p{0.16\textwidth}
  >{\raggedright\arraybackslash}p{0.42\textwidth}
  >{\raggedright\arraybackslash}p{0.40\textwidth}@{}}
\toprule
\textbf{Notebook} & \textbf{Reads} & \textbf{Writes} \\
\midrule
\endhead
A AIS Data Prep & Raw AIS (\texttt{af.get\_ais}), IHS
\texttt{ShipData}/\texttt{tblNameHistory}, \texttt{imf\_port\_boundary.pkl},
\texttt{geoboundary2.pkl} & \texttt{hourly\_ais\_v2/} \\
B IHS Specs & IHS \texttt{ShipData}, \texttt{tblCapacities},
\texttt{tblBoilers}; AE/AB power-factor JSON &
\texttt{ihs\_emissions\_specs\ \{version\}.pkl} / \texttt{.parquet} \\
C Emissions wIHS & \texttt{hourly\_ais\_v2/}, IHS specs parquet,
\texttt{GFW\_effort\_event/\{year\}.csv} &
\texttt{emissions\_data/with\_ihs/} \\
D Emissions woIHS & \texttt{emissions\_data/with\_ihs/} (to fit the
model), \texttt{hourly\_ais\_v2/}, \texttt{GFW\_effort\_event/} &
\texttt{emissions\_data/model/}, \texttt{emissions\_data/no\_ihs/} \\
E Emissions Statistics & \texttt{emissions\_data/with\_ihs/},
\texttt{emissions\_data/no\_ihs/}, \texttt{GFW\_effort\_event/} &
\texttt{emissions\_data/stats/*.csv} \\
(Excel) & \texttt{emissions\_data/stats/*.csv} &
\texttt{Dashboard\ Emissions\ Pacific.xlsx} \\
\bottomrule
\end{longtable}

\subsection{7. Output}\label{output}

The deliverable is \texttt{Dashboard\ Emissions\ Pacific.xlsx}, which
reports estimated emissions --- CO\textsubscript{2},
CO\textsubscript{2}-equivalent, and the air pollutants (SOx, NOx,
CH\textsubscript{4}, CO, N\textsubscript{2}O, NMVOC, PM10, PM2.5, black
carbon) --- over time and by country, vessel type, and fishing gear,
together with vessel counts. It is assembled from the monthly summary
tables produced by notebook~E (the \texttt{emissions\_data/stats/}
outputs). No notebook writes the workbook directly.

\subsection{8. AI Disclosure}\label{ai-disclosure}

Generative AI tools were used to assist in drafting and editing this
documentation, and in formatting the accompanying notebooks. The
notebooks themselves are primarily written by the author. All content
was reviewed by the author, who is responsible for its accuracy.

\end{document}
